\documentclass[10pt,a4paper]{amsart}
\usepackage[margin=19mm]{geometry}
\usepackage{amsmath,amssymb,mathtools}
\usepackage[scr=rsfs,cal=euler]{mathalfa}
\usepackage{xfrac}
\usepackage[unicode,hidelinks,bookmarks=false]{hyperref}
\newcommand{\ie}{i.e.\ }
\newcommand{\cf}{cf.\ }
\newcommand{\resp}{respectively\ }
\newcommand{\Subsection}{\subsection}
\providecommand{\nobreakdash}{\nobreak\hbox{-}}
\setlength{\emergencystretch}{2em}
\multlinegap0pt
\usepackage{bm}
\usepackage{bbm}
\usepackage{dsfont}
\usepackage{xspace}
\usepackage{cancel}
\usepackage{mathtools}
\usepackage{amsthm}
\makeatletter
\@ifundefined{theorem}{\newtheorem{theorem}{Theorem}}{}
\@ifundefined{lemma}{\newtheorem{lemma}[theorem]{Lemma}}{}
\makeatother
\title[Elliptic criticality versus Volterra memory in indirect chem]{Elliptic criticality versus Volterra memory in indirect chemotaxis cascades}
\author{Louis Shuo Wang}
\date{Source: arXiv:2606.12771v1 (CC BY 4.0)}
\begin{document}
\maketitle

\noindent\emph{Source and licence.} Elliptic criticality versus Volterra memory in indirect chemotaxis cascades. Version arXiv:2606.12771v1 (CC BY 4.0), distributed under the Creative Commons Attribution 4.0 International licence, as recorded in the arXiv licence metadata for this exact version and shown on the abstract page https://arxiv.org/abs/2606.12771v1. Full rights notice: licenses/arxiv-2606.12771-CC-BY-4.0.txt.\par\smallskip

\noindent\textit{Editorial scope.} Counted unit: the Lemma whose source label is \texttt{lem:log-kernel-four-dim} in the article (source0.tex lines 486--500, source label \texttt{lem:log-kernel-four-dim}), delivered together with the article's own complete original proof of it (source0.tex lines 502--555, one \texttt{proof} environment of 54 lines). No mathematics is written, repaired or re-derived: statement and proof are the article's own text line for line. Required context retained uncounted: none. Every definition, display and label of those lines is retained unchanged. Omitted from this package and not redistributed: 1--485, 501--501, 556--2434. Nothing inside a retained excerpt is omitted or altered: every retained body is delivered exactly as the article's lines are written, so no omission receipt of any kind accompanies this package. Citations: the retained excerpts carry no citation command at all, so no bibliography entry of the article is transcribed and no reference is redistributed. Reference adaptation: none was needed and none was made; every reference inside the delivered unit resolves inside it, so no unresolved reference or citation is produced. Printed slips kept literal: no printed word, symbol, hypothesis or number of the counted statement or of its proof was altered. Layout and class: the article's own class is \texttt{sn-jnl} with packages including graphicx, multirow, amsmath, amssymb, amsfonts, amsthm, mathrsfs, appendix, xcolor, textcomp, manyfoot, booktabs, algorithm, algorithmicx; the wrapper is the lane's pinned \texttt{amsart} editorial wrapper at 10pt on A4, which restates the theorem-like environments the retained text needs and adds the article's own macro definitions that the retained text needs (none), with extra wrapper packages (bm, bbm, dsfont, xspace, cancel, mathtools). The delivered numbering is the unit's own. Assets: no retained span contains a figure environment, a graphics-inclusion command, a PDF-page inclusion, a table or a TikZ picture; no image file is copied into this package, the package declares no asset dependency and no image file is redistributed with it. Licence: version arXiv:2606.12771v1 of this article is distributed under the Creative Commons Attribution 4.0 International licence, as recorded in the arXiv licence metadata for this exact version and shown on the abstract page https://arxiv.org/abs/2606.12771v1. A full rights notice is delivered at licenses/arxiv-2606.12771-CC-BY-4.0.txt. Characters: every retained excerpt is pure ASCII, so the delivered engine, XeLaTeX with its default Latin Modern text font, reports no missing character.
\begin{lemma}[Four-dimensional logarithmic kernel]
\label{lem:log-kernel-four-dim}
Let \(G_\tau\) denote the local kernel of \(K_\tau\).  In dimension \(N=4\), its
leading near-diagonal singularity is logarithmic:
\[
        G_\tau(x,y)
        =
        \frac{1}{8\pi^2\tau}\log\frac1{|x-y|}
        +
        R_\tau(x,y),
\]
where \(R_\tau\) is less singular than the logarithmic term locally away from
the boundary.  In dimensions \(N\ne4\), the corresponding principal singularity
is of power-law type \(|x-y|^{4-N}\).
\end{lemma}

\begin{proof}
Fix $N=4$ and work locally near the diagonal; boundary and lower-order
contributions are absorbed into $R_\tau$ by standard elliptic parametrix
theory.

Assume first $\tau\neq1$. Partial fractions in the symbol variable give the
operator identity
\[
   K_\tau=(I-\Delta)^{-1}(I-\tau\Delta)^{-1}
        =\frac{1}{1-\tau}(I-\Delta)^{-1}
         -\frac{\tau}{1-\tau}(I-\tau\Delta)^{-1},
\]
which exhibits $K_\tau$ as a difference of two order $-2$ Bessel
potentials. Each summand is individually as singular as $(-\Delta)^{-1}$, so
the asserted order $-4$ behaviour can only arise from a cancellation; we make
it explicit.

Let $G^{(\mu)}$ denote the kernel of $(I-\mu\Delta)^{-1}$, $\mu>0$. Writing
$G^{(\mu)}=\tfrac{1}{4\pi^2\mu}|x|^{-2}+h_\mu$ and using
$-\Delta\bigl(\tfrac{1}{4\pi^2}|x|^{-2}\bigr)=\delta$ in $\mathbb R^4$, one
finds $(I-\mu\Delta)h_\mu=-\tfrac{1}{4\pi^2\mu}|x|^{-2}$. Since
$\Delta\log|x|=2|x|^{-2}$ in $\mathbb R^4$, the leading local solution is
$h_\mu=\tfrac{1}{8\pi^2\mu^2}\log|x|+\mathcal O(1)$. Hence, near the diagonal,
\[
   G^{(1)}(x)=\frac{1}{4\pi^2}|x|^{-2}+\frac{1}{8\pi^2}\log|x|+\mathcal O(1),
   \qquad
   G^{(\tau)}(x)=\frac{1}{4\pi^2\tau}|x|^{-2}
                +\frac{1}{8\pi^2\tau^2}\log|x|+\mathcal O(1).
\]
Substituting into the partial-fraction identity, the $|x|^{-2}$ coefficients
sum to
\[
   \frac{1}{1-\tau}\cdot\frac{1}{4\pi^2}
   -\frac{\tau}{1-\tau}\cdot\frac{1}{4\pi^2\tau}=0,
\]
so the second-order singularities cancel exactly. The surviving logarithmic
coefficient is
\[
   \frac{1}{1-\tau}\cdot\frac{1}{8\pi^2}
   -\frac{\tau}{1-\tau}\cdot\frac{1}{8\pi^2\tau^2}
   =\frac{1}{8\pi^2(1-\tau)}\Bigl(1-\frac1\tau\Bigr)
   =-\frac{1}{8\pi^2\tau}.
\]
Therefore $G_\tau(x,y)=\tfrac{1}{8\pi^2\tau}\log\frac{1}{|x-y|}+R_\tau(x,y)$
with $R_\tau$ bounded near the diagonal, which is the asserted expansion.

The borderline case $\tau=1$ (where the partial fraction degenerates into a
double pole) follows either by continuity in $\tau$ of the coefficient
$\tfrac{1}{8\pi^2\tau}$, or directly: $K_1=(I-\Delta)^{-2}$ has high-frequency
symbol $|\xi|^{-4}$, whose fundamental solution is the biharmonic kernel
$\tfrac{1}{8\pi^2}\log\frac1{|x|}$ in $\mathbb R^4$. In dimensions $N\neq4$ the
same partial-fraction/symbol analysis yields the power-law principal
singularity $|x-y|^{4-N}$.
\end{proof}

\end{document}
